%% -----------------------------------------------------------
\section{Survey Methodology}
%% -----------------------------------------------------------

\subsection{Search Strategy}

This survey employed a systematic search procedure across seven databases: ACM Digital Library, PubMed, PsycINFO, IEEE Xplore, MEDLINE, Scopus, and Google Scholar. The search strategy combined three concept groups using Boolean operators: AI and digital mental health terms (\textit{chatbot, conversational agent, large language model, digital mental health}); marginalized population identifiers (\textit{LGBTQ+, disabled, neurodivergent, autistic, ADHD, culturally diverse, non-Western, racial minority, low-income}); and outcome measures (\textit{depression, anxiety, safety, equity, bias, accessibility}). Additional literature was identified through citation searching of included studies and relevant systematic reviews.

\subsection{Inclusion and Exclusion Criteria}

Studies were included if they: examined AI-powered or digital mental health interventions; included marginalized or vulnerable populations as primary samples or subgroup analyses; reported empirical data from \gls{rct}s, observational studies, mixed-methods research, or systematic reviews; were published in peer-reviewed venues between 2020 and 2025; and addressed effectiveness, safety, accessibility, or equity considerations. Studies were excluded if they focused solely on general populations without marginalized group analysis, were purely theoretical without empirical data, examined non-AI digital interventions, were conference abstracts or editorials, or were non-peer-reviewed grey literature.

\subsection{Corpus Composition}

Thirty-five studies met inclusion criteria. The corpus encompasses five \gls{rct}s~\cite{he_2022,karkosz_2024,liu_2023,meheli_2022,rollwage_2024}; seven systematic reviews and meta-analyses~\cite{abd_alrazaq_2020,casu_2024,dehbozorgi_2025,li_2023,liu_2023,omar_2025,haider_2024}; five observational studies~\cite{de_freitas_2024,drydakis_2022,khan_seo_2025,tanni_2024,papadopoulos_2025}; five qualitative and mixed-methods studies~\cite{aleem_2024,fowler_2023,naderbagi_2024,soubutts_2024,killian_2023}; one implementation science study~\cite{smith_2023}; two policy and institutional reports~\cite{apa_2025,who_2025}; five additional systematic reviews and critical analyses~\cite{lawrence_2024,meadi_2025,pfohl_2024,timmons_2023,robinson_2024}; one bibliometric analysis~\cite{zhang_bibliometric_2025}; and four complementary studies~\cite{huang_2024,palmer_2025,straw_2020,dehghani_2024}.

\subsection{Synthesis Approach}

Following Webster and Watson's~\cite{webster_watson_2002} guidance on concept-centric literature review methodology, analysis employed thematic synthesis organized around the three research questions. Studies were categorized by intervention type and population focus. A concept matrix (Table~\ref{tab:concept_matrix}) organizes the literature by theme, exposing cross-cutting patterns that author-by-author review obscures.

Two seminal works anchor the survey. Khan and Seo~\cite{khan_seo_2025} provide the first comprehensive empirical documentation of \gls{dmh} accessibility failures for blind users, establishing accessibility as a safety issue rather than a design preference. Karkosz et al.~\cite{karkosz_2024} provide one of the first rigorous \gls{rct}s demonstrating both promise and limitations of culturally-adapted AI mental health interventions in a non-English language, establishing an evidence-based methodology for cross-cultural adaptation.
